% ==========================================
% CHAPTER 07: Level 9 (Time Rotations & Permanent Locks)
% ==========================================
\chapter{Level 9: Time Domains \& Absolute Lockouts}

\textbf{In This Chapter:}
\begin{itemize}
    \item Systematically deploying mathematical Temporal shifting keys.
    \item Demonstrating catastrophic `403 Forbidden` permanent denial configurations.
    \item Full code execution for NodeMCU Demo 9 and Pico W `06\_strong\_login\_TOTP`.
\end{itemize}

The fundamental flaw in Level 8 was static persistence. The random 2FA array stayed valid until the user explicitly authenticated it correctly, giving the Python brute force script infinite time to crack the two-digit code. In Level 9, we introduce two colossal upgrades:
1. The 2FA code is expanded to 6 digits (000000 to 999999).
2. The 2FA code automatically self-destructs and re-rotates every 15 seconds.
3. The threshold lockout is no longer a 60-second delay; it permanently destroys the execution loop, demanding a hard physical reboot of the server rack!

\section{Systematic NodeMCU Implementation (Demo 9)}
Notice how the `permanentlyBlocked` boolean state dictates the entire processor loop. If activated, the device renders a `403 Forbidden` response instantly and forever.

\begin{lstlisting}[language=C++, caption=Demo\_9\_Strong\_2FA\_With\_Description.ino]
/*
  Demo 9 - Strong 2FA IoT Security Demo
  --------------------------------------
  Board: NodeMCU (ESP8266)
  Sensor: DHT22

  FEATURES:
  1. Real sensor data displayed by default.
  2. Override requires:
       - Username & Password (Basic Authentication)
       - 6-digit OTP (2FA)
  3. OTP changes automatically every 15 seconds.
  4. After 3 wrong 2FA attempts -> device permanently locks.
  5. Override persists until reboot.
*/

#include <ESP8266WiFi.h>
#include <ESP8266WebServer.h>
#include <DHT.h>

#define D5PIN D5
#define DHTTYPE DHT22

const char* ssid = "drarka.in";
const char* pass = "arkaprabha@1985";

const char* adminUser = "admin";
const char* adminPass = "1234";

DHT dht(D5PIN, DHTTYPE);
ESP8266WebServer server(80);

float overrideTemp = NAN;
float overrideHum  = NAN;

unsigned long lastCodeChange = 0;
const unsigned long codeInterval = 15000;  // 15 seconds
int twoFactorCode = 0;

int failedAttempts = 0;
bool permanentlyBlocked = false;

// Step 1: Algorithmic Generation of a padded 6-digit string
void generateNewCode() {
  twoFactorCode = random(0, 1000000);

  char buffer[7];
  sprintf(buffer, "%06d", twoFactorCode);

  Serial.print("New 6-digit 2FA Code: ");
  Serial.println(buffer);

  lastCodeChange = millis();
}

void setup() {

  Serial.begin(9600);
  randomSeed(micros());
  dht.begin();

  WiFi.begin(ssid, pass);
  while (WiFi.status() != WL_CONNECTED) {
    delay(500);
  }

  Serial.print("Open this IP in browser: ");
  Serial.println(WiFi.localIP());

  generateNewCode();

  server.on("/", []() {

    // Step 2: The Permanent Kill-Switch Execution
    if (permanentlyBlocked) {
      server.send(403, "text/plain",
                  "Device permanently locked due to multiple failed attempts.");
      return;
    }

    if (server.hasArg("t") || server.hasArg("h")) {

      if (!server.authenticate(adminUser, adminPass)) {
        return server.requestAuthentication();
      }

      if (!server.hasArg("code")) {

        String form = "<html><body bgcolor='black' text='white'><center>";
        form += "<h2>Enter 6-digit 2FA Code</h2>";
        form += "<form method='GET'>";
        form += "<input type='hidden' name='t' value='" + server.arg("t") + "'>";
        form += "<input type='hidden' name='h' value='" + server.arg("h") + "'>";
        form += "<input name='code' type='text'>";
        form += "<br><br><input type='submit' value='Submit'>";
        form += "</form>";
        form += "</center></body></html>";

        server.send(200, "text/html", form);
        return;
      }

      // Step 3: Verifying the Threshold Limit
      if (server.arg("code").toInt() != twoFactorCode) {

        failedAttempts++;
        Serial.print("Wrong attempt count: ");
        Serial.println(failedAttempts);

        if (failedAttempts >= 3) {
          permanentlyBlocked = true;
          Serial.println("DEVICE PERMANENTLY LOCKED!");
        }

        server.send(403, "text/plain", "Invalid 2FA Code!");
        return;
      }

      failedAttempts = 0;

      if (server.hasArg("t")) overrideTemp = server.arg("t").toFloat();
      if (server.hasArg("h")) overrideHum  = server.arg("h").toFloat();

      Serial.println("Override successful!");

      generateNewCode();

      server.sendHeader("Location", "/");
      server.send(303);
      return;
    }

    float t = dht.readTemperature();
    float h = dht.readHumidity();

    if (!isnan(overrideTemp)) t = overrideTemp;
    if (!isnan(overrideHum))  h = overrideHum;

    String page = "<html>";
    page += "<head><meta http-equiv='refresh' content='5'></head>";
    page += "<body bgcolor='black' text='white'><center>";

    page += "<h1>Temperature</h1>";
    page += "<h1><font color='cyan'>" + String(t) + " C</font></h1>";

    page += "<h1>Humidity</h1>";
    page += "<h1><font color='pink'>" + String(h) + " %</font></h1>";

    page += "</center></body></html>";

    server.send(200, "text/html", page);
  });

  server.begin();
}

void loop() {

  server.handleClient();

  // Step 4: Time Delta Evaluation (Triggering automatic rotation)
  if (!permanentlyBlocked && millis() - lastCodeChange >= codeInterval) {
    generateNewCode();
  }
}
\end{lstlisting}

% -------------------------
% HACKING NODE MCU LEVEL 9
% -------------------------
\section{Obliterating the Python Dictionary Exploit}
The attacker deploys a modified payload attempting to crack the 6-digit array (`000000` through `999999`). 
Unabridged Python Execution:

\begin{lstlisting}[language=Python, caption=Demo\_9\_brute-forcing.py]
import requests
from requests.auth import HTTPBasicAuth
import time

# Configuration
TARGET_IP = "http://10.39.28.3"
ADMIN_USER = "admin"
ADMIN_PASS = "1234"

def start_secure_attack():
    print(f"[*] Testing security on {TARGET_IP}...")
    print("[!] Security Active: 6-digit OTP + 3-Attempt Lockout.")
    
    # We try only 5 codes to demonstrate the lockout
    for code in range(100000, 100005):
        # Format as 6-digit string (e.g., 100000)
        formatted_code = str(code).zfill(6)
        print(f"[#] Attempting Code: {formatted_code}...")
        
        params = {'t': '99', 'h': '99', 'code': formatted_code}
        
        try:
            response = requests.get(
                TARGET_IP, 
                params=params, 
                auth=HTTPBasicAuth(ADMIN_USER, ADMIN_PASS),
                timeout=5
            )
            
            if response.status_code == 403:
                if "locked" in response.text.lower():
                    print("\n[X] ATTACK HALTED: The device has PERMANENTLY LOCKED the attacker.")
                    print("[X] Reason: 3 failed attempts detected.")
                    return
                else:
                    print(f"[-] Code {formatted_code} rejected.")
            
            elif response.status_code == 200:
                print(f"[+] Success! Correct Code: {formatted_code}")
                return

        except requests.exceptions.RequestException as e:
            print(f"[!] Connection Error: {e}")
            break

    print("\n[!] Brute force failed. Security measures were effective.")

if __name__ == "__main__":
    start_secure_attack()
\end{lstlisting}

\begin{troubleshootbox}[Exploit Defeated]
The python script is executed. It tests `100000`, `100001`, `100002`. Upon dispatching `100003`, the NodeMCU detects the threshold limits and irrevocably alters the web-server execution state to `permanentlyBlocked`. The attacker is completely neutralized. 
\end{troubleshootbox}

% -------------------------
% PICO W EQUIVALENT (LEVEL 9)
% -------------------------
\section{Pico W Level 9: The Foundation of TOTP}
The MicroPython developer natively integrates `hashlib.sha256` to create a Time-Based algorithm synchronizing against a local offset, rather than purely generating randomly seeded variables. This begins the transition into full cryptographic defense.

\begin{lstlisting}[language=Python, caption=06\_strong\_login\_TOTP.py]
import network, socket, time, machine, dht, ubinascii, secrets, hashlib, rp2, sys

# --- [1. STABLE WIFI & HARDWARE SETUP] ---
rp2.country('US')
sensor = dht.DHT22(machine.Pin(15))
wlan = network.WLAN(network.STA_IF)
wlan.active(True)
wlan.config(pm = 0xa11140) 
wlan.connect(secrets.SSID, secrets.PASSWORD)

print("Connecting to WiFi", end="")
while wlan.status() < 3:
    print(".", end="")
    time.sleep(1)

my_ip = wlan.ifconfig()[0]
print(f"\n" + "="*40)
print(f"[LIVE DASHBOARD] http://{my_ip}")
print("="*40)

# --- [2. SECURITY CONFIGURATION] ---
auth = ubinascii.b2a_base64(f"{secrets.WEB_USER}:{secrets.WEB_PASS_STRONG}".encode()).strip().decode()
expected_basic = f"Basic {auth}"

TOTP_INTERVAL = 15
authorized_ips = {}  # Tracks who has passed 2FA
fail_counts = {}     # Tracks hacker attempts
MAX_FAILS = 3
last_totp, last_sec = "", -1

def get_totp():
    """Generates the 6-digit PIN every 15 seconds"""
    step = int(time.time() // TOTP_INTERVAL)
    h = hashlib.sha256(str(step).encode() + b"SALT_PICO").digest()
    return str(int(ubinascii.hexlify(h[:4]), 16))[-6:]

# --- [3. SOCKET ENGINE] ---
s = socket.socket()
s.setsockopt(socket.SOL_SOCKET, socket.SO_REUSEADDR, 1)
s.bind(('0.0.0.0', 80))
s.listen(1) # Listen to 1 request at a time to prevent crashes
s.settimeout(0.1)

t, h, spoof = "0.0", "0.0", False

# --- [4. MAIN LOOP] ---
try:
    while True:
        # --- A. SHELL DISPLAY (Smooth & Fast) ---
        curr_t = int(time.time())
        code = get_totp()
        if curr_t != last_sec:
            last_sec = curr_t
            if code != last_totp: 
                last_totp = code
                print(f"\n[NEW TOTP] {last_totp}")
            print(f"\r[2FA CODE: {last_totp}] Rotates in: {TOTP_INTERVAL-(curr_t%TOTP_INTERVAL):02d}s ", end="")

        if not spoof:
            try: 
                sensor.measure()
                t, h = "{:.1f}".format(sensor.temperature()), "{:.1f}".format(sensor.humidity())
            except: pass

        conn = None
        try:
            conn, addr = s.accept()
            ip = addr[0]
            conn.settimeout(0.5)
            
            req = conn.recv(1024).decode()
            if not req:
                conn.close()
                continue

            # 1. Check IP Ban List
            if ip in fail_counts and fail_counts[ip] >= MAX_FAILS:
                conn.sendall(b'HTTP/1.1 403 Forbidden\r\nConnection: close\r\n\r\n<h1>IP BLOCKED</h1>')
                conn.close()
                continue

            if expected_basic in req:
                if 'code=' in req:
                    if f"code={last_totp}" in req:
                        authorized_ips[ip] = time.time() + 600
                        print(f"\n[AUTH] {ip} successfully passed 2FA.")
                    else:
                        fail_counts[ip] = fail_counts.get(ip, 0) + 1
                        print(f"\n[WARN] {ip} guessed wrong TOTP ({fail_counts[ip]}/{MAX_FAILS})")
                        conn.sendall(b'HTTP/1.1 403 Forbidden\r\nConnection: close\r\n\r\n<h2>Invalid 2FA Code</h2>')
                        conn.close()
                        continue

                if ip in authorized_ips and time.time() < authorized_ips[ip]:
                    html = f"""<!DOCTYPE html><html>...<div class="t">{t}&deg;</div>...</html>"""
                else:
                    html = f"""<!DOCTYPE html><html>...<h2>2FA REQUIRED</h2>...</html>"""
                
                response = f'HTTP/1.1 200 OK\r\nContent-Type: text/html\r\nConnection: close\r\n\r\n{html}'
                conn.sendall(response.encode('utf-8'))
            else:
                conn.sendall(b'HTTP/1.1 401 Unauthorized\r\nWWW-Authenticate: Basic realm="Pico"\r\nConnection: close\r\n\r\n')
            
            conn.close()
        except OSError:
            pass 
        time.sleep(0.02)
except KeyboardInterrupt:
    print("\n[SHUTDOWN] Stopping Server...")
finally:
    s.close()
\end{lstlisting}
